\section{Conclusion}

Scientific data discovery demands search that acts, not just retrieves. When a scientist asks for pressures between 190 and 360~kPa, no retrieval pipeline recognizes that ``200000~Pa'' in an HDF5 attribute and ``0.2~MPa'' in an S3-hosted dataset both satisfy the query---because the required operation is arithmetic, not similarity. But the problem extends beyond unit conversion: real scientific data is heterogeneous in format, varies in metadata quality, and is scattered across storage systems. At the scale of 10 billion blobs, brute-force inspection is impossible. We addressed this gap with CLIO Search, a framework for agentic scientific discovery that separates \emph{reasoning about what to search} from \emph{executing the search}.

The system implements a four-stage pipeline. The \emph{Discover} stage profiles each namespace to learn what data exists and what metadata is available (completing in under 12ms). The \emph{Reason} stage selects retrieval branches based on corpus characteristics, achieving 95.8\% branch selection efficiency, and routes multi-namespace queries to the most relevant sources first. The \emph{Transform \& Search} stage executes science-aware operators---dimensional conversion and formula normalization---as deterministic, pluggable retrieval branches alongside BM25 and dense vector search. The \emph{Execute} stage provides federated data access across six storage backends with multi-hop refinement.

The results validate the approach across controlled benchmarks, real observatory data, a live external catalog, and distributed HPC infrastructure. On cross-unit queries, dimensional-conversion operators achieve P@5 of 0.99 where the strongest text baseline scores 0.40; on 288 real NOAA observatory documents, precision improves by 87\%. The three-path agentic comparison shows that wrapping CLIO behind a single tool call reduces LLM token consumption by 52\% and wall time by 91\% compared to a native agent baseline, while answering all 10 queries. Cross-unit correctness is verified at distributed scale: on 70K arXiv documents across 4 DeltaAI GH200 nodes, querying in Pa, kPa, bar, and atm returns identical hit counts, with sub-50ms latency. IOWarp integration demonstrates that CLIO can serve as the search layer for tiered blob storage, inspecting 0.5\% of the corpus per query where the native BlobQuery API returns everything unfiltered. A deterministic no-LLM fallback achieves 96\% of LLM-assisted quality, confirming that the operators---not the language model---are the primary source of gains.

Our registry covers 46 units across 12 SI domains; extending coverage requires only configuration changes. Compound units and free-text unit parsing remain unsupported. Distributed experiments on NCSA DeltaAI validate scaling to 2.5M documents across 4~GH200 nodes with 0.94 weak scaling efficiency; IOWarp integration is validated at 10K blobs. Scaling to petabyte-scale parallel filesystems with billions of blobs remains future work.

Several directions follow naturally. Integrating NDP-MCP~\cite{mcp_hpc2025} as a Discover data source would enable the agent to learn what datasets exist across facilities. Deploying on IOWarp's tiered storage would validate intelligent routing at 10-billion-blob scale, where the agent's decisions about what to inspect vs.\ skip become critical. Replacing heuristic fusion with learned reranking~\cite{agenticr2026} would improve the agentic loop. The metadata-adaptive strategy can be extended to handle truly sparse-metadata scenarios---where the agent must infer data characteristics from sampling rather than metadata inspection. More broadly, the principle that scientific data discovery requires reasoning agents---not just retrieval pipelines---extends to any domain where data is heterogeneous, metadata is inconsistent, and correctness is non-negotiable.
